\section{Discussion}

The kill test reveals that compensation plays two distinct roles in transformer grokking. It accelerates learning in favourable regimes and it steers the system toward stable generalizing circuits. When compensation is blocked learning either slows dramatically or converges to brittle solutions that do not generalize.

The strong seed 0 result matches the original picture of Le Chatelier style homeostasis. At $\omega = 0.5$ the perturbed head is weakened and other heads pick up the slack. This redistribution of variance introduces instability into the memorization plateau and accelerates escape. When those compensatory heads are frozen the system can no longer rebalance. It remains stuck near the perturbed configuration and must solve parity by brute force search through weight space. The result is a $4.9$ times slowdown.

Seeds 1 and 3 show that this is not an artefact of a single random initialization. When compensation is blocked the cost can be as small as a twenty percent delay, but in no seed does freezing compensation make learning faster while preserving generalization. The system benefits from being able to redistribute variance across heads even when the magnitude of that benefit is modest.

Seed 2 is more revealing. Here freezing compensatory heads causes the model to grok quickly but badly. $T_{\text{grok}}$ is low and the final test accuracy collapses toward chance. The free run groks much later but reaches high test accuracy. This pattern suggests that homeostatic compensation is not just a speed up mechanism. It constrains the space of solutions the model explores. With compensation active the system is pushed toward circuits that maintain some conserved quantity such as total information flow or representation rank. With compensation blocked the system can satisfy the loss with a narrow shortcut that does not respect that conserved structure.

Viewed through a thermodynamic lens the suppressor heads and their companions define an equilibrium manifold in weight space \cite{kringelbach2024thermodynamics}. Training near this manifold exhibits critical slowing down. Small perturbations decay slowly and the system spends many steps in a metastable memorizing state. Strong perturbations push the system off the manifold. Compensation then acts to pull it back while simultaneously destabilizing the memorization plateau. When compensation is blocked the manifold is broken. The system either crawls along a distorted plateau for many more steps or hops into a different attractor that solves the loss without restoring equilibrium.
